\section{The problem}

Bard, Robertson, and Sorace made a methodological claim that was larger than a
single elicitation format. Acceptability judgments, they argued, could be
treated as graded measurements rather than as informal binary verdicts
\citep{BardEtAl1996}. That move helped make experimental syntax a
measurement enterprise: judgments could be collected from naive participants,
scaled, compared across groups, and subjected to ordinary statistical tests.

Whether magnitude estimation was necessary for that advance is a harder
question.
Magnitude estimation promised a way around the scale restrictions of
traditional judgment tasks, giving participants an open response scale and
preserving fine distinctions. But later acceptability work didn't settle into
one privileged elicitation method. Likert ratings, forced choice, yes/no tasks,
and Thurstonian models all became part of the measurement toolkit
\citep{sprouse2013,LangsfordEtAl2018,sprouse2017}.

This paper separates two claims that are often allowed to travel together. One
claim is strong and still plausible: acceptability is graded, measurable, and
methodologically serious. The other is weaker: magnitude estimation has a
practical measurement advantage over simpler formal tasks. The first claim can
survive even if the second doesn't.

This empirical question isn't new. The methodology literature has already
compared elicitation formats and largely deflated magnitude
estimation's special status. \textcite{schutze2016} reports that most or all of
its advertised advantages over Likert tasks have since been refuted, pointing to
\textcite{weskott2011} and \textcite{sprouse2013}; later work weighed the formats
on convergence \citep{sprouse2013}, statistical power \citep{sprouse2017},
and reliability and response style \citep{LangsfordEtAl2018}.

This paper instead contributes a clearer routing of the evidence. Existing
method-comparison work already weakens magnitude estimation's special status;
the reassessment asks what the
surviving formal-judgment signal licenses us to project. In the approved
reanalysis of the 2013 and 2017 comparison data and a post-outcome robustness
multiverse, formal judgment
formats preserve stable item-level contrast ordering across the tested
contrasts, while magnitude estimation adds no robust separate practical
ordering there.

Seen this way, the reassessment doesn't reject Bard et al.'s programme. It
tests what should be credited to magnitude estimation itself, and what
should instead be credited to the broader formalization of acceptability
judgment methodology.
